% Chapter 27: Grammar Detective: The Great Error Hunt

\chapter{Grammar Detective: The Great Error Hunt}

\section{Pedagogical Purpose}
This chapter integrates everything already learned—agreement, tense, punctuation, word choice, and sentence construction—into single passages containing several different kinds of error at once, and gives learners a repeatable checklist for finding them all, not just the one they happen to notice first.

\begin{learningobjectives}
The learner can:
\begin{itemize}
    \item name the five broad categories used in the self-editing checklist.
    \item apply the checklist to find more than one kind of error in the same passage.
    \item identify which earlier chapter's rule explains a given error.
    \item correct an error and explain why it was wrong, not just supply the fix.
    \item edit a full passage for agreement, tense, punctuation, word choice, and sentence construction together.
\end{itemize}
\end{learningobjectives}

\section{Prerequisite Check}
\begin{quickcheck}
Name the chapter (by topic, not number) that taught you the rule needed to fix each sentence.
\begin{enumerate}
    \item "The bunch of flowers are beautiful."
    \item "Yesterday, we go to the market."
    \item "wheres my pencil box"
\end{enumerate}
\end{quickcheck}

\begin{rememberbox}
You have now learned rules for agreement, tense, punctuation, word choice, and sentence building across many separate chapters. This chapter does not teach a new rule—it teaches you how to use everything you already know \textbf{at the same time}, the way a real editor checks a whole piece of writing.
\end{rememberbox}

\section{Chapter Opener}
Leo the parrot has "written" a diary entry and is very proud of it, but Ms. Sharma spots trouble immediately.

\textbf{Leo's Diary Entry:} "Yesterday I fly to the window and I sees a big bird. Its wings was huge! The bird land on a branch and it dont move for a long time. I is very excited."

\textbf{Rohan:} "Leo, I count at least... five mistakes in there!"

\textbf{Maya:} "Some are tense mistakes, some are agreement mistakes, and I think 'its' should be 'it's' somewhere too..."

\textbf{Ms. Sharma:} "Exactly! Leo's diary is the perfect case for a Grammar Detective who has learned everything in this book. Today, instead of hunting for just one kind of mistake, you'll hunt for \textbf{all} of them—using a checklist real editors use."

\section{Language Experience}
\textbf{Leo's Diary Entry, Before and After}

\begin{quote}
\begin{tcolorbox}[colback=gray!5, colframe=gray!50, sharp corners, boxsep=0pt, left=5pt, right=5pt, top=5pt, bottom=5pt, breakable]
\textbf{Before:} Yesterday I fly to the window and I sees a big bird. Its wings was huge! The bird land on a branch and it dont move for a long time. I is very excited.
\end{quote}
\end{tcolorbox}

\begin{quote}
\begin{tcolorbox}[colback=gray!5, colframe=gray!50, sharp corners, boxsep=0pt, left=5pt, right=5pt, top=5pt, bottom=5pt, breakable]
\textbf{After:} Yesterday I flew to the window and I saw a big bird. Its wings were huge! The bird landed on a branch and it didn't move for a long time. I was very excited.
\end{quote}
\end{tcolorbox}

\section{Look and Notice}
\begin{thinkbox}
\begin{enumerate}
    \item How many separate corrections were needed to fix Leo's diary entry?
    \item Can you sort the corrections into groups—which ones are about \textit{tense}, and which are about \textit{agreement}?
    \item Was there a punctuation fix hiding in there too ("dont" $\rightarrow$ "didn't")?
    \item If you had only checked for tense errors, what might you have missed?
\end{enumerate}
\end{thinkbox}

\section{Discover / Explicit Teaching}
A real editor does not check a piece of writing only once, looking for only one thing. Instead, they work through a \textbf{checklist}, checking for one category of error at a time, all the way through the passage, before moving to the next category.

\section{Grammar Word}
\begin{grammarword}{SELF-EDITING CHECKLIST}
A \textbf{self-editing checklist} is a fixed list of categories a writer checks, one at a time, all the way through a piece of writing, to catch as many different kinds of error as possible—rather than reading once and fixing only whatever is noticed first.

\textit{Examples:}
\begin{itemize}
    \item Checking a whole paragraph for agreement errors first, then going back through the same paragraph again just for tense, then again just for punctuation.
\end{itemize}

\textit{Non-example:} Reading a passage once and fixing only the first mistake spotted, then assuming the rest must be correct, is not using a checklist—it usually misses errors from categories the writer wasn't actively looking for.

\textit{Quick Check:} Why might checking a passage five separate times (once per category) catch more errors than checking it only once?
\end{grammarword}

\section{Core Explanation}
\subsection*{Agreement errors (Chapter 13)}
Example: "The basket of apples \textbf{are} heavy." $\rightarrow$ should be \textbf{is}.

\subsection*{Tense errors (Chapters 14-15)}
Example: "Yesterday, we \textbf{go} to the market." $\rightarrow$ should be \textbf{went}.

\subsection*{Punctuation errors (Chapter 16)}
Example: "wheres my pencil box" $\rightarrow$ should be "\textbf{W}here's my pencil box\textbf{?}"

\subsection*{Word choice errors (Chapters 18-19)}
Example: "\textbf{Its} raining today." $\rightarrow$ should be \textbf{It's}.

\subsection*{Sentence construction errors (Chapter 20)}
Example: "The dog. It barked. It ran." $\rightarrow$ "The dog barked and ran."

\section{Worked Example}
\begin{examplebox}
\textbf{Passage to edit:} "My sister and me goes to the park every Sunday. We plays cricket and eat lunch. Last week it rain and we gets wet. Its very funny now when we remember it."

\textit{Step 1 (Agreement):} Check each subject-verb pair. "My sister and me goes" $\rightarrow$ the subject "My sister and I" is plural, so it should be "go," not "goes." Also, "and me" should be "and I" as the subject of the sentence.
\textit{Step 2 (Tense):} "Last week it rain" signals the past but uses the base form—should be "rained." "We gets wet" should be "got wet" to match the past-tense story.
\textit{Step 3 (Punctuation):} "Its very funny" should be "\textbf{It's} very funny" (it is).
\textit{Step 4 (Word choice):} No homophone or vague-word errors remain once the apostrophe above is fixed.
\textit{Step 5 (Sentence construction):} Each sentence is already complete and in natural order, so no further changes are needed here.

\textbf{Answer:} "My sister and I \textbf{go} to the park every Sunday. We \textbf{play} cricket and eat lunch. Last week it \textbf{rained} and we \textbf{got} wet. \textbf{It's} very funny now when we remember it."
\end{examplebox}

\section{Guided Practice}
\begin{activitybox}{Activity A: Identification}
Sort each error into one of the five checklist categories (agreement, tense, punctuation, word choice, sentence construction).
\begin{enumerate}
    \item "The dogs barks loudly."
    \item "Yesterday we go to the fair."
    \item "wheres my bag"
\end{enumerate}
\end{activitybox}

\section{Independent Practice}
\begin{activitybox}{Independent Practice}
\begin{enumerate}
    \item Edit this passage using the full five-step checklist: "My uncle live in Delhi. He visit us last month and he bring gifts for everyone. There was a big cake and we all has some."
    \item Choose one sentence you corrected in Question 1 and explain, in your own words, which chapter's rule it broke.
    \item Write a two-sentence passage of your own containing exactly two different kinds of error (from two different categories) for a partner to find.
    \item \textit{Reasoning:} Why might it be harder to spot a tense error in a sentence that also contains a punctuation error, compared to a sentence with only one mistake?
\end{enumerate}
\end{activitybox}

\section{Summary}
\begin{summarybox}
\begin{itemize}
    \item A \textbf{self-editing checklist} helps a writer check a passage systematically, one category at a time.
    \item The five checklist steps are: \textbf{agreement}, \textbf{tense}, \textbf{punctuation}, \textbf{word choice}, and \textbf{sentence construction}.
    \item Multi-category passages, closer to real writing, often contain more than one kind of error at once.
    \item Not every unusual sentence is an error—some are deliberate \textbf{stylistic choices}.
    \item Explaining \textit{why} an error is wrong is as important as fixing it.
\end{itemize}
\end{summarybox}

\section{Key Terms}
\begin{table}[h!]
\centering
\begin{tabular}{|l|p{9cm}|}
\hline
\textbf{Term} & \textbf{Simple Meaning} \\
\hline
Self-editing checklist & A fixed list of categories checked one at a time when editing writing \\
Multi-category passage & A passage containing more than one kind of grammar error \\
Stylistic choice & A deliberate writing choice that is valid, not a genuine error \\
Agreement & Subject and verb matching in number (Chapter 13) \\
Consistency & Keeping tense (or another feature) steady throughout a passage (Chapter 15) \\
\hline
\end{tabular}
\end{table}

\begin{selfcheck}
I can:
\begin{itemize}
    \item name the five categories in the self-editing checklist.
    \item apply the checklist to find more than one kind of error in a passage.
    \item identify which earlier chapter's rule explains a given error.
    \item correct an error and explain why it was wrong.
    \item edit a full passage using all five checklist categories.
\end{itemize}
\end{selfcheck}